\documentclass{article}
\usepackage{amsmath, amssymb, booktabs, hyperref}

\title{Derivation of Pushing-Medium Calibration Constants:\\
$\mu_G$, $k_{\mathrm{TT}}$, and $k_{\mathrm{Fizeau}}$}
\author{Pushing-Medium Project}
\date{March 2026}

\begin{document}
\maketitle

\begin{abstract}
Three numerical constants appear throughout the Pushing-Medium (PM) numerical
test suite and comparison code: the gravitational index coefficient
$\mu_G = 1.4861 \times 10^{-27}$\,m$^2$\,kg$^{-1}$, the transverse-traceless
power normalisation $k_{\mathrm{TT}} = 1.0$, and the Fizeau drag coupling
$k_{\mathrm{Fizeau}} = 1.0$.  This document records exactly how each value is
obtained, the underlying physical reasoning, and how each value connects to
known General Relativity (GR) results.  The numbers are reproducible from
first principles via the calibration fixture in
\texttt{legacy/Pushing-Medium/tests/fixtures/calibration.py}.
\end{abstract}

% ------------------------------------------------------------------
\section{Background: The PM Index Field for Gravity}
% ------------------------------------------------------------------

In the Pushing-Medium framework, spacetime is modelled as a medium whose
refractive index $n$ encodes the gravitational field.  For a point mass $M$ the
index at distance $r$ is

\begin{equation}
  n(r) = 1 + \frac{\mu_G\, M}{r},
  \label{eq:index}
\end{equation}

where $\mu_G$ is the gravitational index coefficient (a coupling constant with
SI units m$^2$\,kg$^{-1}$).  Light rays are governed by Fermat's principle in
this medium: the deflection angle accumulated along a straight reference path at
impact parameter $b$ is

\begin{equation}
  \alpha_{\mathrm{num}}(\mu_G)
  = \int_{-z_{\max}}^{+z_{\max}}
      \frac{\partial \ln n}{\partial x}\bigg|_{x=b,\,y=0}\, dz
  = \int_{-z_{\max}}^{+z_{\max}}
      \frac{-\mu_G M\, b}{\,n(r)\, r^{3}\,}\, dz,
  \label{eq:alpha_num}
\end{equation}

where $r = \sqrt{b^2 + z^2}$.  The analytic weak-field GR (Schwarzschild)
result that this must reproduce is

\begin{equation}
  \alpha_{\mathrm{GR}} = \frac{4GM}{c^2 b}.
  \label{eq:gr_deflection}
\end{equation}

% ------------------------------------------------------------------
\section{Calibration of $\mu_G$}
\label{sec:mu}
% ------------------------------------------------------------------

\subsection{Theoretical expectation}

Comparing the leading-order expansion of~\eqref{eq:alpha_num} (with $n \approx 1$
in the denominator) to the GR result~\eqref{eq:gr_deflection} immediately gives

\begin{equation}
  \alpha_{\mathrm{num}} \approx
  \mu_G M \int_{-\infty}^{+\infty} \frac{b}{\,(b^2+z^2)^{3/2}\,}\, dz
  = \mu_G M \cdot \frac{2}{b}.
\end{equation}

Setting this equal to~\eqref{eq:gr_deflection} yields

\begin{equation}
  \mu_G^{\mathrm{theory}} = \frac{2G}{c^2}
  = \frac{2 \times 6.67430\times10^{-11}}{(2.99792458\times10^8)^2}
  \approx 1.4862 \times 10^{-27}\;\text{m}^2\,\text{kg}^{-1}.
  \label{eq:mu_theory}
\end{equation}

\subsection{Numerical determination procedure}

Because the full integrand in~\eqref{eq:alpha_num} retains the factor
$1/n(r)$ rather than approximating $n \approx 1$, the exact numerical value
of $\mu_G$ is obtained by a bisection search rather than the closed-form
expression~\eqref{eq:mu_theory}.  The procedure is:

\begin{enumerate}
  \item Fix a solar test case: lens mass $M_\odot = 1.989\times10^{30}$\,kg,
        impact parameter $b_\odot = 6.96\times10^8$\,m (solar radius).
  \item Compute the GR target:
        $\alpha^* = 4GM_\odot/(c^2 b_\odot) \approx 8.481\times10^{-6}$\,rad
        ($\approx 1.75''$).
  \item Bracket $\mu_G$ in the interval
        $[0.25 \times 2G/c^2,\; 4.0 \times 2G/c^2]$.
  \item Run 36 bisection iterations, each evaluating
        $\alpha_{\mathrm{num}}(\mu_G)$ numerically via
        \texttt{index\_deflection\_numeric} with
        $z_{\max} = 200\,b_\odot$, $N = 2500$ steps.
  \item Return the midpoint of the final bracket.
\end{enumerate}

The result converges to

\begin{equation}
  \mu_G = 1.4861356300677034 \times 10^{-27}\;\text{m}^2\,\text{kg}^{-1}.
  \label{eq:mu_result}
\end{equation}

The relative difference from the analytic estimate~\eqref{eq:mu_theory} is
$< 0.01\%$; the small discrepancy arises from the $1/n$ correction in the
integrand and the finite integration window.

\subsection{Physical interpretation}

$\mu_G = 2G/c^2$ is the Schwarzschild radius per unit mass divided by $c$.
Equivalently, it is the coefficient that makes the PM index gradient reproduce
the standard PPN $\gamma = 1$ gravitational lensing result.  Its appearance at
exactly $2G/c^2$ is not accidental: it is the value required so that
Fermat's principle in the medium~\eqref{eq:index} is equivalent, to leading
order, to null geodesics in the linearised Schwarzschild metric.

% ------------------------------------------------------------------
\section{Calibration of $k_{\mathrm{TT}}$}
\label{sec:kTT}
% ------------------------------------------------------------------

\subsection{Role of $k_{\mathrm{TT}}$}

The PM theory predicts gravitational-wave energy emission through a
transverse-traceless (TT) quadrupole formula.  The PM expression for the
radiated power from a circular binary (masses $M_1$, $M_2$, separation $a$) is

\begin{equation}
  P_{\mathrm{PM}} = \frac{32}{5}\,\frac{G^4\,(M_1 M_2)^2(M_1+M_2)}{c^5\,a^5}.
  \label{eq:pm_power}
\end{equation}

The GR (Peters 1964) formula for the same system is

\begin{equation}
  P_{\mathrm{GR}} = \frac{32}{5}\,\frac{G^4\,(M_1 M_2)^2(M_1+M_2)}{c^5\,a^5}.
  \label{eq:gr_power}
\end{equation}

The normalisation factor is defined as

\begin{equation}
  k_{\mathrm{TT}} = \frac{P_{\mathrm{GR}}}{P_{\mathrm{PM}}}.
  \label{eq:kTT_def}
\end{equation}

\subsection{Numerical determination}

The ratio is evaluated at a fiducial double neutron-star-like system:
$M_1 = 1.4\,M_\odot$, $M_2 = 1.3\,M_\odot$, $a = 10^9$\,m.
Because Equations~\eqref{eq:pm_power} and~\eqref{eq:gr_power} share identical
functional forms—the PM TT sector was constructed to reproduce GR
gravitational-wave predictions—the ratio is

\begin{equation}
  k_{\mathrm{TT}} = 1.0.
  \label{eq:kTT_result}
\end{equation}

\subsection{Physical interpretation}

$k_{\mathrm{TT}} = 1$ confirms that the PM quadrupole power formula is
numerically identical to the GR Peters formula at the order implemented.
This is by design: the TT projector and the $32/5$ pre-factor are imported
directly from the linearised GR derivation.  A value differing from unity
would indicate an inconsistency in the PM TT sector that would need to be
resolved before the theory could be used for gravitational-wave predictions.

% ------------------------------------------------------------------
\section{Calibration of $k_{\mathrm{Fizeau}}$}
\label{sec:kFizeau}
% ------------------------------------------------------------------

\subsection{Physical context}

When a gravitational lens moves transversely with speed $v \ll c$, the
deflection angle acquires a first-order correction relative to the static
result.  In GR (with the standard PPN parameter $\gamma = 1$) the
correction is

\begin{equation}
  \alpha_{\mathrm{moving}} \approx \alpha_{\mathrm{static}}
  \left(1 + \frac{v}{c}\right) + \mathcal{O}\!\left(\frac{v^2}{c^2}\right).
  \label{eq:gr_moving}
\end{equation}

In the PM framework this correction is parametrised as

\begin{equation}
  \alpha_{\mathrm{PM,moving}} = \alpha_{\mathrm{static}}
  \left(1 + k_{\mathrm{Fizeau}}\,\frac{v}{c}\right),
  \label{eq:pm_moving}
\end{equation}

where $k_{\mathrm{Fizeau}}$ is the Fizeau drag coupling.

\subsection{Determination}

Matching~\eqref{eq:pm_moving} to the GR result~\eqref{eq:gr_moving} gives
directly

\begin{equation}
  k_{\mathrm{Fizeau}} = \frac{\left(\dfrac{\alpha_{\mathrm{moving}}}
                                          {\alpha_{\mathrm{static}}} - 1\right)}
                              {\dfrac{v}{c}}.
  \label{eq:kFizeau_def}
\end{equation}

With the GR target ratio $\alpha_{\mathrm{moving}}/\alpha_{\mathrm{static}}
= 1 + v/c$, this evaluates to

\begin{equation}
  k_{\mathrm{Fizeau}}
  = \frac{(1 + v/c) - 1}{v/c}
  = 1.0,
  \label{eq:kFizeau_result}
\end{equation}

independent of the chosen test velocity (fiducial value $v = 30$\,km\,s$^{-1}$
is used in the code).

\subsection{Physical interpretation}

$k_{\mathrm{Fizeau}} = 1$ means the PM medium-drag correction to moving-lens
deflection equals the GR first-order $v/c$ correction when $\gamma = 1$.
Physically, the Fizeau drag of light by a moving medium is the non-relativistic
analogue of the gravitomagnetic contribution to bending in GR.  The unit value
is consistent with the relativistic aberration / frame-dragging picture and
acts as a sanity check that the PM kinematic sector correctly reproduces the
PPN $\gamma = 1$ moving-lens result.  

\textit{Note:} the designation ``Fizeau'' refers to the 19th-century Fizeau
drag experiment (partial entrainment of light by a moving medium), and by
analogy to the partial frame-dragging of null rays by a moving gravitational
lens in the PM index picture.

% ------------------------------------------------------------------
\section{Summary}
% ------------------------------------------------------------------

\begin{center}
\begin{tabular}{llll}
\toprule
\textbf{Constant} & \textbf{Value} & \textbf{Method} & \textbf{Reference} \\
\midrule
$\mu_G$ & $1.4861356 \times 10^{-27}$\,m$^2$\,kg$^{-1}$
         & Bisection fit: PM path-integral $\to$ $4GM/c^2b$
         & Sec.~\ref{sec:mu} \\
$k_{\mathrm{TT}}$ & $1.0$
         & Ratio: PM quadrupole power / GR Peters power
         & Sec.~\ref{sec:kTT} \\
$k_{\mathrm{Fizeau}}$ & $1.0$
         & Algebraic: PM moving-lens model $\to$ GR $v/c$ correction
         & Sec.~\ref{sec:kFizeau} \\
\bottomrule
\end{tabular}
\end{center}

\bigskip
All three constants reduce, up to numerical precision, to values expected from
linearised GR with $\gamma = 1$.  This provides a self-consistency check: the
PM framework recovers standard weak-field GR predictions in the appropriate
limits when its index coupling is fixed to $\mu_G = 2G/c^2$.

\bigskip\noindent
\textit{Code reference:}
\texttt{legacy/Pushing-Medium/tests/fixtures/calibration.py},
function \texttt{load\_or\_calibrate()}.

\end{document}
